% Unit 035 English translation of chapter3.tex lines 622--699.
% Source: Methods of Algebra, Volume 2, commit
% 9a5803ff77dd3257484cb177f851a73770a59dd3 (CC BY 4.0).
% stable-unit-id: o014.aljabr2.chapter3.opposite-category-complexes
% Coverage: Section 3.4 in full.

% segment-id: o014.li2.chapter3.opposite-category-complexes.h001
\section{Complexes on the Opposite Category}\label{sec:opposite-cplx}
\hypertarget{o014.aljabr2.chapter3.opposite-category-complexes}{ }

% segment-id: o014.li2.chapter3.opposite-category-complexes.p001
Let $\mathcal{A}$ be an additive category. This section relates complexes on
$\mathcal{A}$ and on $\mathcal{A}^{\opp}$. Since functors involving
$\mathcal{A}^{\opp}$, such as the $\Hom$ functor, occur frequently, this
procedure is simple but necessary; some of the signs involved also require a
little care. Beginning readers are advised to skip this section or merely
skim it on a first reading.

% segment-id: o014.li2.chapter3.opposite-category-complexes.p002
The discussion has three aspects: complexes, homotopy, and mapping cones. The
derived-category version considered later
(Proposition~\sourcecrossref{prop:derived-cat-op}{in the later discussion of
derived categories}) is a direct application of these results. Recall that if
$f:X\to Y$ is a morphism in $\mathcal{A}$, then $f^{\opp}$ denotes the
corresponding morphism $Y\to X$ in $\mathcal{A}^{\opp}$.

% segment-id: o014.li2.chapter3.opposite-category-complexes.dp001
\begin{definition-proposition}\label{def:sigma}
	\index[sym1]{sigma@$\sigma$}
	The isomorphism of additive categories
	$\sigma:\cate{C}(\mathcal{A}^{\opp})\to\cate{C}(\mathcal{A})^{\opp}$
	is defined as follows. For an object $X$ of
	$\cate{C}(\mathcal{A}^{\opp})$, define in $\mathcal{A}$
% segment-id: o014.li2.chapter3.opposite-category-complexes.q001
	\[ (\sigma X)^n := X^{-n}, \quad d_{\sigma X}^n = (-1)^{n+1} \left[ d_X^{-n-1, \opp}: (\sigma X)^n \to (\sigma X)^{n+1} \right], \quad n \in \Z. \]
	For a morphism $f=\left(f^n:X^n\to Y^n\right)_n$ in
	$\cate{C}(\mathcal{A}^{\opp})$, take its image $\sigma f$ to be the
	following morphism in $\cate{C}(\mathcal{A})$:
% segment-id: o014.li2.chapter3.opposite-category-complexes.q002
	\[ (\sigma f)^n := \left(f^{-n} \right)^{\opp} : (\sigma Y)^n \to (\sigma X)^n, \]
	that is, a morphism $\sigma X\to\sigma Y$ in
	$\cate{C}(\mathcal{A})^{\opp}$. For every $m\in\Z$, there is a natural
	isomorphism
% segment-id: o014.li2.chapter3.opposite-category-complexes.q003
	\[ s_m: \sigma \circ [m] \rightiso [-m] \circ \sigma, \]
	where $[-m]$ is regarded as a functor from
	$\cate{C}(\mathcal{A})^{\opp}$ to itself (strictly speaking, it should be
	written $[-m]^{\opp}$).
\end{definition-proposition}
% segment-id: o014.li2.chapter3.opposite-category-complexes.pf001
\begin{proof}
	The general case of $s_m$ reduces step by step to the case $m=1$. Define
	$s_1=(s_{1,X})_X$ as follows. For every $n$, take
% segment-id: o014.li2.chapter3.opposite-category-complexes.g001
	\begin{equation}\label{eqn:sigma-s}\begin{tikzcd}[row sep=tiny]
		s_{1,X}^n : \sigma(X{[1]})^n = X^{1-n} \arrow[r, "\sim"', "{(-1)^{n-1}}" inner sep=0.6em] & X^{1-n} = (\sigma X){[-1]}^n .
	\end{tikzcd}\end{equation}
	The entire verification reduces to
	$s_{1, X}^{n+1} d^n_{\sigma(X[1])} = -d_{X[1]}^{-n-1, \opp} =
	d_X^{-n, \opp} = (-1)^n d_{\sigma X}^{n-1} =
	d^n_{(\sigma X)[-1]} s_{1,X}^n$.
\end{proof}

% segment-id: o014.li2.chapter3.opposite-category-complexes.p003
If the construction $\sigma$ with coefficient $(-1)^{n+1}$ is applied in
both directions, $\sigma^2X$ is naturally isomorphic to $X$ through the
isomorphism whose degree-$n$ component is $(-1)^n\identity_{X^n}$. To obtain
an inverse, define the reverse construction by the same formulas on objects
and morphisms, but use the coefficient $(-1)^n$ on the differential. The
composites in both orders are then exactly the identity functor, so $\sigma$
is indeed an isomorphism of categories.%
\footnote{Translator's note (O014-C036): arguing from
$n(n+1)\equiv0\pmod{2}$, the source states that applying $\sigma$ twice
returns $X$ strictly. Using the same formula in both directions instead gives
the differential $-d_X$, not $d_X$. The signed natural isomorphism and the
corrected inverse construction are recorded above.}

% segment-id: o014.li2.chapter3.opposite-category-complexes.r001
\begin{remark}\label{rem:sigma-vs-cohomology}
	If $\mathcal{A}$ is an abelian category, the signs on
	$d_{\sigma X}^\bullet$ do not change the cohomology introduced in
	Definition~\ref{def:cohomology}; thus
	$\Hm^{-n}(\sigma X)\in\Obj(\mathcal{A})$ corresponds to
	$\Hm^n(X)\in\Obj(\mathcal{A}^{\opp})$.
\end{remark}

% segment-id: o014.li2.chapter3.opposite-category-complexes.p004
We next examine the relation between $\sigma$ and homotopy.

% segment-id: o014.li2.chapter3.opposite-category-complexes.pr001
\begin{proposition}\label{prop:sigma-homotopy}
	The functor $\sigma$ defined above induces an equivalence of additive
	categories
	$\cate{K}(\mathcal{A}^{\opp})\rightiso\cate{K}(\mathcal{A})^{\opp}$.
\end{proposition}
% segment-id: o014.li2.chapter3.opposite-category-complexes.pf002
\begin{proof}
	Let $X,Y$ be objects of $\cate{C}(\mathcal{A}^{\opp})$. For
	$f=(f^k)_k\in\Hom^{-1}_{\cate{C}(\mathcal{A}^{\opp})}(X,Y)$, use the
	following table to define the corresponding
	$\sigma f\in\Hom^{-1}_{\cate{C}(\mathcal{A})}(\sigma Y,\sigma X)$ in
	$\mathcal{A}$:
% segment-id: o014.li2.chapter3.opposite-category-complexes.q004
	\begingroup\setlength{\arraycolsep}{4pt}
	\[\begin{array}{|c|c|c|c|} \hline
		\text{Category} & \multicolumn{3}{c|}{\text{Morphism}} \\ \hline
		\mathcal{A}^{\opp} & X^k \xrightarrow{(-1)^{k+1} f^k} Y^{k-1} & d_Y^{k-1} f^k + f^{k+1} d_X^k & (d^{-1}f)^k \\
		\mathcal{A} & (\sigma Y)^{-k+1} \xrightarrow{(\sigma f)^{-k+1}} (\sigma X)^{-k} & (\sigma f)^{-k+1} d_{\sigma Y}^{-k} + d_{\sigma X}^{-k-1} (\sigma f)^{-k} & d^{-1} (\sigma f)^{-k} \\ \hline
	\end{array}\]
	\endgroup
	This suffices to show
	\[\begin{aligned}
		\Hom_{\cate{K}(\mathcal{A}^{\opp})}(X, Y)
		&\rightiso \Hom_{\cate{K}(\mathcal{A})}(\sigma Y, \sigma X) \\
		&= \Hom_{\cate{K}(\mathcal{A})^{\opp}}(\sigma X, \sigma Y).
	\end{aligned}\]
\end{proof}

% segment-id: o014.li2.chapter3.opposite-category-complexes.p005
Finally, using $\sigma$ and the family of isomorphisms $(s_m)_m$ constructed
above, we compare mapping cones in $\cate{C}(\mathcal{A})$ and
$\cate{C}(\mathcal{A}^{\opp})$. The details are somewhat routine.

% segment-id: o014.li2.chapter3.opposite-category-complexes.pr002
\begin{proposition}\label{prop:sigma-triangulated}
	\index[sym1]{sigma}
	For every morphism $f:X\to Y$ in $\cate{C}(\mathcal{A}^{\opp})$, there
	is a commutative diagram in $\cate{C}(\mathcal{A})$%
	\footnote{Here $\alpha(\cdot)$ and $\beta(\cdot)$ are both defined
	relative to $\cate{C}(\mathcal{A})$. If the diagram is viewed in
	$\cate{C}(\mathcal{A})^{\opp}$, all its arrows must be reversed.}:
% segment-id: o014.li2.chapter3.opposite-category-complexes.g002
	\[\begin{tikzcd}[column sep=large]
		\sigma\left( X{[1]} \right) \arrow[r, "{\sigma(\beta(f))}"] \arrow[d, "{s_{1,X}}"'] & \sigma\left(\Cone(f)\right) \arrow[r, "{\sigma(\alpha(f))}"] \arrow[d, "\theta"] & \sigma Y \arrow[r, "\sigma(f)"] \arrow[d, "\identity"] & \sigma X \arrow[d, "{\identity}"] \\
		(\sigma X){[-1]} \arrow[r, "{\alpha(\sigma(f))[-1]}"' inner sep=0.6em] & \Cone(\sigma(f)){[-1]} \arrow[r, "{\beta(\sigma(f))[-1]}"' inner sep=0.6em] & \sigma Y \arrow[r, "{\sigma(f)}"' inner sep=0.6em] & \sigma X
	\end{tikzcd}\]
	where $\theta$ is a canonical isomorphism and $s_{1,X}$ is the
	isomorphism supplied by Definition--Proposition~\ref{def:sigma}.
\end{proposition}
% segment-id: o014.li2.chapter3.opposite-category-complexes.pf003
\begin{proof}
	If $\Cone(f)$ were replaced by $X[1]\oplus Y$, the statement would be
	easy and the required isomorphism would simply be
	$(s_{1,X},\identity_{\sigma Y})$. The complication is that the matrix
	$d_{\Cone(f)}$ has the off-diagonal entry $f[1]$. Nevertheless, we follow
	the same method: using \eqref{eqn:sigma-s}, for every $n\in\Z$ define the
	isomorphism
% segment-id: o014.li2.chapter3.opposite-category-complexes.q005
	\begin{multline*}
		\theta^n: \sigma\left(\Cone(f)\right)^n = \sigma\left(X[1]\right)^n \oplus (\sigma Y)^n \\
		\xrightarrow[\sim]{(s_{1,X}, \identity)^n = ((-1)^{n-1} \identity, \identity) } (\sigma X)[-1]^n \oplus (\sigma Y)^n \xlongequal{\text{interchange}} \left( (\sigma Y) \oplus (\sigma X)[-1] \right)^n .
	\end{multline*}
	Under this identification, the morphism $d_{\sigma(\Cone(f))}^n$ in
	$\mathcal{A}$ corresponds to
% segment-id: o014.li2.chapter3.opposite-category-complexes.q006
	\[ \begin{pmatrix}
		d_{\sigma Y}^n & 0 \\
		* & d_{(\sigma X)[-1]}^n
	\end{pmatrix} : \left( (\sigma Y) \oplus (\sigma X)[-1] \right)^n \to \left( (\sigma Y) \oplus (\sigma X)[-1] \right)^{n+1} . \]
	As noted at the start of the proof, the diagonal entries pose no problem;
	we need only determine the lower-left entry. By
	Definition--Proposition~\ref{def:sigma}, this entry is $(-1)^{n+1}$ times
	the following composite in $\mathcal{A}$:
% segment-id: o014.li2.chapter3.opposite-category-complexes.g003
	\[\begin{tikzcd}[row sep=tiny]
		(\sigma Y)^n \arrow[r] & (\sigma(X[1]))^{n+1} \arrow[r, "\sim"] & (\sigma X)[-1]^{n+1} \\
		Y^{-n} \arrow[equal, u] \arrow[r, "{(f^{-n})^{\opp}}"'] & X^{-n} \arrow[equal, u] \arrow[r, "{s_{1, X}^{n+1} = (-1)^n}"'] & X^{-n} \arrow[equal, u]
	\end{tikzcd}\]
	Thus the result is $-\sigma(f)^n$. It follows that
% segment-id: o014.li2.chapter3.opposite-category-complexes.q007
	\[ \theta := (\theta^n)_{n \in \Z}: \sigma(\Cone(f)) \rightiso \Cone(\sigma(f))[-1] \]
	is indeed an isomorphism in $\cate{C}(\mathcal{A})$. Once this is
	established, verifying commutativity presents no essential difficulty.
\end{proof}

% segment-id: o014.li2.chapter3.opposite-category-complexes.p006
Of course, the results of this section also apply to chain complexes and
extend to the case in which $\mathcal{A}$ is a $\Bbbk$-linear category.
